This section motivates the survey, defines the evaluation crisis facing frontier Large Language Models (LLMs), and summarises the seven evolving threads that the rest of the paper unifies. We open with the model-release discontinuity, then sketch the saturation arc of static benchmarks, and finally preview the multi-axis taxonomy that organises the remainder. The release of GPT-3 in 2020 and ChatGPT in November 2022 marked a discontinuity in how language models are built, deployed, and judged {[}Brown et al., 2020; Ouyang et al., 2022{]}. Where earlier natural-language-processing systems were assessed through narrow task accuracies on the General Language Understanding Evaluation (GLUE) benchmark and its harder successor SuperGLUE {[}Wang et al., 2018; Wang et al., 2019{]}, modern decoder-only transformers with hundreds of billions of parameters are expected to perform competently across hundreds of heterogeneous tasks: graduate-level science question answering on GPQA-Diamond {[}Rein et al., 2024{]}, competition-level mathematical reasoning on MATH {[}Hendrycks et al., 2021{]}, real-world software engineering on SWE-bench {[}Jimenez et al., 2024{]}, multi-turn dialogue on MT-Bench {[}Zheng et al., 2023{]}, long-context retrieval on RULER up to 128k tokens {[}Hsieh et al., 2024{]}, and agentic tool use on GAIA {[}Mialon et al., 2024{]}. The intellectual challenge that this survey addresses is to organise a now-vast and rapidly fragmenting evaluation literature into a coherent map that researchers, practi\ldots{}
